\begin{abstract}
This paper reports the design, object-oriented firmware architecture, and
hardware-validated performance of Steistat, our own low-cost,
open-hardware potentiostat built around the Analog Devices AD5940
electrochemical analog front end (AFE) and a Seeed Studio XIAO RP2040
microcontroller, built directly from the AD5940 datasheet, standard
three-electrode potentiostat design practice, and the AD5940
evaluation-kit reference firmware. Every electrochemical-method class is
designed under one polymorphic base class with a genuine inheritance
relationship, and the central parameter store, \texttt{C\_DataStorage},
uses private state with dedicated accessor methods throughout. Measurement
latency is bounded and datasheet-matched: each conversion settles within
$\approx$\SI{13.3}{\milli\second} ($\approx$\SI{75}{\hertz} ceiling), with
every wait loop timeout-protected. Cyclic voltammetry (CV) against a traceable
Randles dummy cell (nominal DC resistance \SI{10560}{\ohm}) recovers
$R=\SI{10605.1\pm0.9}{\ohm}$ across five replicate sweeps
($+0.427\,\%$ error, $R^2\geq0.999998$, $\%\mathrm{RSD}=0.0084\,\%$), and a
companion check against a plain \SI{10}{\kilo\ohm} resistor recovers
\SI{10005.9}{\ohm} ($+0.06\,\%$ error, $R^2=0.9999988$). This accuracy and
latency followed a direct, register-level verification of our firmware
against Analog Devices' own AD5940 reference examples, which identified
and corrected concrete defects -- spanning the excitation/sensing
signal path, the ADC conversion convention, the on-board
calibration-resistor declaration, and several further register- and
timing-level mismatches -- with the complete hardware design files, bill
of materials, and firmware source openly available for independent
reproduction. CA's initial, step-onset current is correctly signed and
monotonic in the applied potential across five step levels ($-100$ to
$+200$~mV), and SWV/DPV produce smooth, bounded differential-current
curves with no railing or discontinuities. We also report, honestly and
without correction in this paper, a real, currently open limitation
surfaced during this same diagnostic work: the \emph{sustained},
steady-state current reported by CA, SWV, and DPV does not track the
applied potential, a defect distinct from those fixed here and, on
the evidence available, localized upstream of the sensing and conversion
stage rather than in it. The paper positions its contribution around a
diagnosed, evidence-based, hardware-validated measurement chain, rather
than an unsubstantiated claim that every implemented method is presently
accurate.
\end{abstract}

\begin{keyword}
Potentiostat \sep AD5940 \sep object-oriented design \sep embedded systems
\sep electrochemical instrumentation \sep firmware debugging \sep
reference-code verification \sep dummy cell \sep open-source hardware
\end{keyword}
